% RQ5 Conclusion: Inter-DAO Cooperation

\subsection{Summary}

We have presented a systematic analysis of inter-DAO cooperation through multi-agent simulation. Across \experimentruns{} simulation runs testing five cooperation scenarios, we characterized the conditions under which DAOs can effectively collaborate and the mechanisms that enable stable partnerships.

Key findings:

\begin{enumerate}
    \item \textbf{Success-surplus tradeoff}: Homogeneous DAOs cooperate easily but generate modest value. Heterogeneous coalitions are harder to form but more valuable.

    \item \textbf{Fairness enables sustainability}: Explicit fairness mechanisms (Shapley division) are essential for stable cooperation between unequal partners.

    \item \textbf{Hub structures improve coordination}: Central coordinators reduce negotiation costs and enable larger multi-party coalitions.

    \item \textbf{Overlapping membership facilitates trust}: Cross-DAO token holders create natural alignment channels that increase cooperation probability.

    \item \textbf{Competition doesn't preclude cooperation}: Even rival DAOs can collaborate on clear public goods that benefit all parties.
\end{enumerate}

\subsection{Contributions}

This paper contributes:

\begin{enumerate}
    \item \textbf{Cooperation taxonomy}: Classification of inter-DAO relationship structures
    \item \textbf{Success conditions}: Identification of factors enabling stable cooperation
    \item \textbf{Mechanism design guidance}: Recommendations for joint proposal structures and fairness mechanisms
    \item \textbf{Multi-DAO simulation framework}: Infrastructure for studying cross-organization dynamics
\end{enumerate}

\subsection{Practical Recommendations}

For DAOs considering cooperation:

\begin{enumerate}
    \item \textbf{Identify complementarities}: Seek partners with different strengths, not duplicates of yourself
    \item \textbf{Build relationships first}: Trust-building precedes successful formal cooperation
    \item \textbf{Design for fairness}: Explicit division rules prevent disputes and enable asymmetric partnerships
    \item \textbf{Create overlap}: Encourage cross-DAO participation through shared membership or delegation
    \item \textbf{Consider coordination structures}: For multi-party cooperation, designate or rotate coordinator roles
\end{enumerate}

\subsection{Implications}

For the DAO ecosystem, our results suggest that inter-DAO cooperation is not just possible but potentially highly valuable---if properly structured. The billions held in DAO treasuries could fund ecosystem-wide public goods if cooperation mechanisms mature.

For researchers, our framework enables systematic study of multi-organization dynamics in decentralized contexts. The intersection of coalition theory, mechanism design, and blockchain governance offers rich terrain for future investigation.

\subsection{Closing Remarks}

DAOs face challenges that no single organization can address alone. Public goods, standards, security---these require coordination across organizational boundaries. Our simulations demonstrate that such coordination is achievable, but requires intentional mechanism design.

The future of decentralized governance likely involves not just individual DAOs but networks of cooperating organizations. Understanding the dynamics of inter-DAO cooperation is essential to building that future.

We release our simulation framework to enable the community to extend this research and develop the practical infrastructure for DAO-to-DAO relations.
